%%%%%%%%%%%%
% Preamble %
%%%%%%%%%%%%

\documentclass
  [11pt, a4paper, english, openright, sumario = tradicional, twoside, brazil]
  {abntex2}

%%%%%%%%%%%%%
%% Packages %
%%%%%%%%%%%%%

%%%%%%%%%%%%%%
%%% Encoding %
%%%%%%%%%%%%%%

\usepackage[utf8]{inputenc}
\usepackage[T1]{fontenc}

%%%%%%%%%%%%%%%%%%%%%
%%% ABNT Compliance %
%%%%%%%%%%%%%%%%%%%%%

\usepackage{indentfirst}
\usepackage{microtype}

%%%%%%%%%%%%%%
%%% Glossary %
%%%%%%%%%%%%%%

\let\glossitem\undefined
\usepackage[portuguese]{gloss}

%%%%%%%%%%%%%%%%%%%
%%% Miscellaneous %
%%%%%%%%%%%%%%%%%%%

\usepackage{amsmath}
\usepackage[backend = biber, indent, style = abnt]{biblatex}
\usepackage{color}
\usepackage{csquotes}
\usepackage{graphicx}
\usepackage{palatino}
\usepackage{paralist}
\usepackage{xspace}

%%%%%%%%%%%%%%%%%%
%% Handy Aliases %
%%%%%%%%%%%%%%%%%%

\newcommand{\treesitter}{\textit{Tree-Sitter}\xspace}
\newcommand{\witchcooking}{\textit{Witch Cooking}\xspace}

%%%%%%%%%%%%%%%%%%
%% Configuration %
%%%%%%%%%%%%%%%%%%

%%%%%%%%%%%%%%%%%%%%%%%%%%%%%%%%%%%%%%%%%
%%% ABNT's NBR 14724 Metrics Compliance %
%%%%%%%%%%%%%%%%%%%%%%%%%%%%%%%%%%%%%%%%%

\setlength{\parindent}{1.3cm}
\setlength{\parskip}{0.2cm}

%%%%%%%%%%%%%
%%% abntex2 %
%%%%%%%%%%%%%

\renewcommand{\ABNTEXchapterfont}{\fontseries{b}}

%%%%%%%%%%%%%%%%%%%%%%%%%%%%%%
%%%% Cover & Title Page Info %
%%%%%%%%%%%%%%%%%%%%%%%%%%%%%%

\titulo
  {\witchcooking: Formatação de Código via Linguagem de Consulta do
   \treesitter}
\autor{Átila Gama Silva}
\local{Irecê}
\data{\the\year}
\orientador{Rafael Xavier}
\instituicao
  {Instituto Federal de Educação, Ciência e Tecnologia da Bahia – Campus Irecê}
\tipotrabalho{Trabalho de Conclusão de Curso}
\preambulo
  {\textit{Trabalho de Conclusão de Curso} apresentado ao
   \textit
     {Curso Técnico em Análise e Desenvolvimento de Sistemas do Instituto
      Federal de Educação, Ciência e Tecnologia da Bahia – Campus Irecê},
   como requisito parcial para obtenção do diploma de
   \textit{Técnico em Análise e Desenvolvimento de Sistemas}, mediante a
   orientação do professor \imprimirorientador.}

%%%%%%%%%%%%%%%%%%%
%%%% IFBA's Cover %
%%%%%%%%%%%%%%%%%%%

\renewcommand
  {\imprimircapa}
  {\begin{capa}
     \ABNTEXchapterfont
     \ABNTEXsectionfontsize
     \center
     \includegraphics[width=\textwidth]{ifba-irece-logo.png}
     \par\vfill\imprimirautor
     \par\vfill{\ABNTEXchapterfontsize\imprimirtitulo}
     \vfill\vfill
     \par\imprimirlocal\break\imprimirdata
   \end{capa}}

%%%%%%%%%%%%%%%%%%%%%%%%
%%%% IFBA's Title Page %
%%%%%%%%%%%%%%%%%%%%%%%%

\makeatletter
\renewcommand
  {\folhaderostocontent}
  {\begin{center}
     \ABNTEXchapterfont
     \ABNTEXsectionfontsize
     \imprimirautor
     \par\vfill\vfill{\ABNTEXchapterfontsize\imprimirtitulo}
     \vfill
     {\normalsize
      \mdseries
      \hspace{.45\textwidth}
      \begin{minipage}{.5\textwidth}
        \SingleSpacing
        \imprimirpreambulo
      \end{minipage}}
     \vfill\vfill
     \imprimirlocal\break\imprimirdata
   \end{center}}
\makeatother

%%%%%%%%%%%%%%
%%%% Summary %
%%%%%%%%%%%%%%

\setlength{\cftlastnumwidth}{\cftsectionnumwidth}
\cftsetindents{part}{0em}{\cftlastnumwidth}
\cftsetindents{chapter}{0em}{\cftlastnumwidth}
\cftsetindents{section}{0em}{\cftlastnumwidth}
\cftsetindents{subsection}{0em}{\cftlastnumwidth}
\cftsetindents{subsubsection}{0em}{\cftlastnumwidth}
\cftsetindents{paragraph}{0em}{\cftlastnumwidth}
\cftsetindents{subparagraph}{0em}{\cftlastnumwidth}

\renewcommand
  {\tocprintchapter}
  {\addtocontents{toc}{\cftsetindents{chapter}{0em}{\cftlastnumwidth}}}

\renewcommand
  {\tocinnonumchapter}
  {\addtocontents{toc}{\cftsetindents{chapter}{\cftlastnumwidth}{0em}}
   \cftinserthook{toc}{A}}

\renewcommand
  {\tocprintchapternonum}
  {\addtocontents{toc}{\cftsetindents{chapter}{\cftlastnumwidth}{0em}}}

\renewcommand
  {\tocpartapendices}
  {\addtocontents{toc}{\cftsetindents{part}{\cftlastnumwidth}{0em}}
   \cftinserthook{toc}{A}}

\renewcommand
  {\tocpartanexos}
  {\addtocontents{toc}{\cftsetindents{part}{\cftlastnumwidth}{0em}}
   \cftinserthook{toc}{A}}

\setlength{\cftbeforechapterskip}{1em}

%%%%%%%%%%%%%%%%%%
%%% Bibliography %
%%%%%%%%%%%%%%%%%%

\addbibresource{references.bib}

%%%%%%%%%%%%%%
%%% Glossary %
%%%%%%%%%%%%%%

\setglosslabel{\bfseries#2\ifglossshort{ (#3)}{}}
\renewenvironment{thegloss}
  {\chapter*{\glossname}
   \markboth{\glossname}{\glossname}
   \addcontentsline{toc}{chapter}{\glossname}}
  {\stopglosslist}
\makegloss

%%%%%%%%%%%%%%
%%% hyperref %
%%%%%%%%%%%%%%

\def\chapterautorefname{capítulo}
\def\Itemautorefname{item}

%%%%%%%%%%%%%%
%%% Paralist %
%%%%%%%%%%%%%%

\setdefaultenum{(I)}{(A)}{}{}

%%%%%%%%%
%%% PDF %
%%%%%%%%%

\definecolor{blue}{RGB}{41,5,195}
\makeatletter
\hypersetup
  {pdftitle       = {\@title},
   pdfauthor      = {\@author},
   pdfsubject     = {\witchcooking: Formatação de Código via Linguagem de
                     Consulta do \treesitter},
   pdfcreator     = {\LaTeX with \abnTeX},
   pdfkeywords    = {formatação de código}
                    {\textit{prettyprint}}
                    {formatação personalizada}
                    {\treesitter}
                    {linguagem de consulta},
   colorlinks     = true,
   linkcolor      = blue,
   citecolor      = blue,
   filecolor      = magenta,
   urlcolor       = blue,
   bookmarksdepth = 4}
\makeatother

%%%%%%%%%%%%
% Document %
%%%%%%%%%%%%

\begin{document}
  \selectlanguage{brazil}
  \frenchspacing

  %%%%%%%%%%%%%%
  % Pretextual %
  %%%%%%%%%%%%%%

  \imprimircapa
  \imprimirfolhaderosto

  %%%%%%%%%%%%
  % Abstract %
  %%%%%%%%%%%%

  \setlength{\absparsep}{18pt}

  \begin{resumo}
    A formatação de código é fundamental no desenvolvimento de software,
    permitindo estabelecer aspectos desejados como a padronização e
    legibilidade do código, que impactam positivamente o ciclo de vida do
    software. Em ambientes de desenvolvimento modernos, é comum o uso de
    ferramentas para automatizar a formatação de código. No entanto, essas
    ferramentas geralmente apresentam limitações no número de linguagens
    suportadas, nas opções de configuração e não permitem a definição de
    estilos de formatação personalizados. Visando superar ou reduzir as
    limitações frequentemente encontradas nas ferramentas convencionais de
    formatação, neste trabalho é apresentado o \witchcooking, um software
    prototípico que tem como objetivos
    \begin{inparaenum}
      \item abranger uma gama de linguagens de programação; além de
      \item permitir que o usuário defina seus próprios estilos de formatação.
    \end{inparaenum}

    \noindent
    \textbf{Palavras-chave}:
    formatação de código;
    \textit{prettyprint};
    formatação personalizada;
    \treesitter;
    linguagem de consulta.
  \end{resumo}

  %%%%%%%%%%%%%%%%%%%%%
  % Table of Contents %
  %%%%%%%%%%%%%%%%%%%%%

  \pdfbookmark[0]{\contentsname}{toc}
  \tableofcontents*
  \cleardoublepage

  %%%%%%%%%%%%%%%%%%%%
  % Textual Elements %
  %%%%%%%%%%%%%%%%%%%%

  \textual

  %%%%%%%%%%%%%%%%%
  %% Introduction %
  %%%%%%%%%%%%%%%%%

  \chapter{Introdução}

  A flexibilidade presente na sintaxe de linguagens de programação permite que
  diferentes arranjos dum mesmo código compartilhem um valor sintático
  equivalente. Essa característica possibilita a formatação do código de acordo
  com aspectos desejados, como a legibilidade, que é fundamental no ciclo de
  vida do software
  \cites[546]{buse-2009-learning}[1]{oliveira-2020-evaluating}.
  Consequentemente, a formatação é frequentemente utilizada para estabelecer um
  nível satisfatório de legibilidade em bases de código.

  A formatação de código pode ser realizada manualmente pelo programador,
  embora esse processo possa ser demorado, especialmente em grandes bases de
  código, exigindo do programador um tempo que poderia ser empregado em outra
  tarefa. Além disso, a formatação manual também pode ser falha e inconsistente
  devido à suscetibilidade humana ao erro, podendo ser agravada quando há
  múltiplos programadores em uma base de código.

  Para evitar os problemas inerentes da formatação manual, é comum a utilização
  de softwares que automatizam a formatação de forma determinística, tornando o
  código padronizado e consistente. No entanto, é comum que essas ferramentas
  de formatação sejam limitadas a uma linguagem ou a uma família de linguagens
  de programação. Além disso, algumas delas são opinativas, assim restringindo
  as possiblidades de personalização pelo usuário. Finalmente, é também comum
  que essas ferramentas de formatação não permitam que o usuário defina estilos
  de formatação personalizados.

  Em ambientes de desenvolvimento modernos, além da utilização de ferramentas
  que automatizam a formatação de código, é também comum a utilização do
  \treesitter \cite{tree-sitter-2018-tree}, um sistema de análise sintática de
  código aberto que foi disponibilizado ao público geral no GitHub
  primeiramente em 2019, tendo sido inicialmente desenvolvido por Max
  Brunsfeld. Desde seu lançamento, o \treesitter tem ganhado popularidade na
  comunidade de desenvolvedores devido dentre outros motivos à
  \begin{inparaenum}
    \item sua capacidade de suportar várias linguagens de programação; além de
    \item sua linguagem de consulta, a qual permite a realização de buscas
          complexas na árvore analisada de um código fonte.
  \end{inparaenum}

  Visando superar ou reduzir as limitações previamente mencionadas, quais são
  frequentemente presentes nas ferramentas convencionais de formatação, este
  trabalho tem como objetivo geral desenvolver um software prototípico para a
  formatação de código. Nomeada antecipadamente de \witchcooking por
  \textcite{silva-2023-witch}, a ferramenta a ser desenvolvida tem como
  objetivos
  \begin{inparaenum}
    \item abranger uma gama de linguagens de programação; além de
    \item permitir que o usuário defina seus próprios estilos de formatação.
  \end{inparaenum}

  Para alcançar o objetivo geral proposto, este trabalho tem como objetivos
  específicos
  \begin{inparaenum}
    \item apresentar e contextualizar ferramentas de formatação de código
          conceituadas, além de abordar seus métodos de formatação;
    \item conceituar, contextualizar e abordar as tecnologias do sistema
          \treesitter; e, por fim,
    \item utilizar o \treesitter para desenvolver e atender às aspirações do
          \witchcooking.
  \end{inparaenum}

  %%%%%%%%%%%%%%%%%%%%%%%%%%%%%%%
  %% The Source Code Formatting %
  %%%%%%%%%%%%%%%%%%%%%%%%%%%%%%%

  \chapter{A Formatação de Código-Fonte}

  Desde os primórdios da computação, métodos foram desenvolvidos para garantir
  que a saída impressa fosse formatada de maneira esteticamente agradável
  \apud{harris-1956-keyboard}[1]{yelland-2015-rfmt}. Esses métodos ganharam
  popularidade sob o termo \textit{``\gloss{prettyprinting}''}, que se refere à
  formatação visual de diversos tipos de conteúdo. No desenvolvimento de
  software, o \textit{prettyprinting} é conhecido como formatação de código,
  uma prática histórica e comum que envolve a adoção de convenções estilísticas
  para estruturar o código-fonte. Existem diversas terminologias utilizadas
  para se referir às ferramentas que realizam a formatação do código-fonte,
  incluindo
  \begin{inparaenum}
    \item ``formatadores de código'', ou \textit{``code formatters''} em
          inglês;
    \item \textit{``prettyprinters''}; ou ainda
    \item \textit{``beautifiers''}.
  \end{inparaenum}

  Durante as décadas de 60 e 70, a linguagem de programação LISP%
  \footnote
    {O termo ``LISP'' (acrônimo para
     \textit{``\textbf{LIS}t \textbf{P}rocessing''})
     inicialmente se referia à linguagem de programação desenvolvida por
     \textcite{mccarthy-1960-recursive}, porém, com a disseminação de dialetos
     da linguagem, o termo também passou a ser utilizado para se referir à
     família de linguagens derivadas da original.}
  proporcionou condições favoráveis para o avanço da formatação de código
  \cite[2]{yelland-2015-new}.
  \begin{inparaenum}
    \item LISP apresentava uma sintaxe distinta e expressiva, baseada em listas
          e estruturas aninhadas delimitadas por parênteses%
          \footnote
            {Na notação original de McCarthy, eram utilizadas
             \textit{M-expressions} entre colchetes para representar
             expressões. Essas \textit{M-expressions} seriam posteriormente
             traduzidas em \textit{\gloss[Word]{sexp}s}. Por exemplo, a
             \textit{M-expression} \boxed{\texttt{car[cons[A; B]]}} seria
             equivalente à \textit{S-expression}
             \boxed{\texttt{(car (cons A B))}}.
             No entanto, assim que a linguagem LISP foi implementada, os
             programadores prontamente adotaram o uso das
             \textit{S-expressions} em vez das \textit{M-expressions}, tornando
             as \textit{S-expressions} a forma predominante de representar a
             estrutura do código na linguagem.},
          o que tornava o código-fonte praticamente ilegível e incompreensível
          caso não fosse devidamente organizado. Além disso,
    \item LISP era \gloss{homoiconic}, ou seja, permitia a manipulação do
          código via dados da própria linguagem.
  \end{inparaenum}
  Essas duas características intrínsecas do LISP fundamentaram o surgimento de
  \begin{inparaenum}
    \item \textit{prettyprinters} para a linguagem, visando aperfeiçoar a
          legibilidade do código-fonte escrito ou emitido; além de
    \item novas técnicas e abordagens para a formatação de código em geral.
  \end{inparaenum}
  Esses avanços contribuíram significativamente para o conhecimento e o
  aprimoramento de práticas de formatação de código em trabalhos subsequentes.

  Em 1967, Bill Gosper desenvolveu o \textit{GRINDEF} (acrônimo para
  \textit
    {``\textbf{GRIND}%
       \footnote
         {O termo \textit{``grind''} era utilizado em alguns círculos de LISP
          como sinônimo para \textit{prettyprinters}.}
       \textbf{F}unction''}%
  ), considerado o primeiro \textit{prettyprinter} a mensurar o tamanho das
  linhas e ter ciência de sua localização no arquivo
  \cites{gosper-2023-twubblesome}{griesemer-2022-cultural}. Essa ferramenta
  implementava o algoritmo \textit{recursive re-predictor}, que, como descrito
  por \textcite{goldstein-1973-pretty}, percorria as árvores de nós
  representando as listas e imprimia os nós considerando seus tamanhos e a
  quantidade máxima de caracteres que ainda poderiam ser inseridos em uma
  linha.

  Posteriormente, \textcite{hearn-1979-one} propuseram um algoritmo mais
  elaborado para \textit{prettyprinting} que utilizava um par de corrotinas%
  \footnote
    {Na obra, os autores enfatizam que a implementação com o uso de corrotinas
     é tão simples que poderia ser simulada -- sem grandes dificuldades -- em
     linguagens que não possuíam suporte nativo para corrotinas
     \cite[53]{hearn-1979-one}.},
  uma
  \begin{inparaenum}
    \item responsável por produzir uma sequência de caracteres que
          representavam o programa sendo impresso; e outra
    \item responsável por decidir como esses caracteres seriam exibidos.
  \end{inparaenum}
  Essas corrotinas se comunicavam por meio de um buffer \gloss{fifo},
  permitindo que as decisões de formatação fossem adiadas até que houvesse
  informações suficientes disponíveis para tomá-las com confiabilidade.

  No ano seguinte, \textcite{oppen-1980-prettyprinting} apresentou em sua
  influente obra, intitulada \textit{``Prettyprinting''}, um algoritmo inovador
  de formatação de código-fonte. Embora apresentasse semelhanças ao algoritmo
  proposto por \textcite{hearn-1979-one}, destacava-se por sua capacidade de
  formatar código derivado de qualquer linguagem de programação. Para realizar
  essa tarefa, o algoritmo necessitava que o código-fonte fosse anotado com
  espaços em branco%
  \footnote
    {Quebras de linha, avanços de formulário e alimentações de linha também
     eram tratados como espaços em branco pelo algoritmo.}
  e delimitadores especiais para marcar o início e fim de blocos logicamente
  contíguos. Assim, o código a ser fornecido ao algoritmo precisaria ser
  processado por uma ferramenta intermediária capaz de compreender a sintaxe da
  linguagem e fornecer um código anotado de forma adequada, permitindo que o
  algoritmo realizasse a formatação apropriada.

  \dots % FIXME

  Recentemente, tem se tornado cada vez mais evidente a crescente relevância
  dos softwares opinativos \cite{eccles-2015-rise}, também conhecidos em inglês
  como \textit{\gloss{opinionated} softwares}. Esses softwares preestabelecem
  práticas estritivas, visando deliberadamente frustrar ou dificultar práticas
  destoantes \cite{lance-2021-about}. Similarmente, os \textit{prettyprinters}
  primitivos, tais como o SOAP \cite{scowen-1971-soap}, eram naturalmente
  restritivos devido à simplicidade e às limitações dos algoritmos de
  formatação existentes. Conforme mencionado anteriormente, à medida que os
  algoritmos de formatação de código-fonte eram aprimorados, tornando-se mais
  robustos e sofisticados, os \textit{prettyprinters} passaram a proporcionar
  mais personalizações aos usuários. No entanto, atualmente, a abordagem
  opinativa também tem sido adotada por softwares modernos de formatação de
  código-fonte \cites{black-2023-black}[5--6,8]{griesemer-2022-cultural}. Esses
  softwares geralmente têm como objetivos
  \begin{inparaenum}
    \item padronizar e
    \item garantir a consistência do código; além
    \item eliminar o tempo e esforços necessários para escolher os estilos
          ideais de formatação; e, finalmente, consequentemente e
          principalmente
    \item aprimorar o ambiente de desenvolvimento colaborativo.
  \end{inparaenum}

  Um exemplo proeminente de \textit{prettyprinter} opinativo é o
  \textit{\textbf{Black}} \cite{black-2023-black}, um formatador de código
  Python, cuja logomarca parodia a da fabricante automobilística Ford, e cujo
  irônico slogan \textit{\textbf{``any color you like''}} faz alusão à famosa
  frase
  \textit
    {``any customer can have a car painted \textbf{any colour that he wants so
     long as it is black}''}
  \cite[p. 72, grifo nosso]{ford-1922-my}, dita por Henry Ford, o fundador da
  fabricante. Assim como Ford adotou uma abordagem de produção em massa na
  fabricação dos \textit{Model Ts} com uma única opção de cor para aumentar a
  eficiência \cite{ford-1922-my}, o \textit{Black} segue um princípio análogo.
  Ao impor um estilo estrito de formatação, minimizando o espaço para
  personalização ou preferências individuais, o \textit{Black} oferece
  velocidade, determinismo, além de economia de tempo e esforços para questões
  mais importantes \cite{black-2023-black}.

  Sumarizando, desenvolver um \textit{prettyprinter} não é uma tarefa simples
  \cite[55]{hughes-1995-design}. Durante o processo de formatação, é necessário
  \begin{inparaenum}
    \item ter conhecimento das estruturas gramaticais da linguagem específica
          em questão. Além disso, o objetivo do \textit{prettyprinter} é
    \item \label{inparaenum:format} formatar o código de maneira otimizada e
          agradável para o leitor \cite[55]{hughes-1995-design}, levando em
          consideração as melhores práticas de formatação. Em alguns casos,
          também é importante
    \item permitir que o usuário possa optar por estilos de formatação
          que atenda às suas preferências individuais. Adicionalmente, as
          sintaxes das linguagens de programação costumam ser flexíveis em
          relação à estruturação do código-fonte. Assim, para atender o
          \autoref{inparaenum:format}, geralmente é necessário
    \item utilizar algoritmos que calculem o layout mais adequado de acordo
          com o contexto
          \cites
            [12]{jasper-2023-clang}
            {goldstein-1973-pretty}
            {hughes-1995-design}
            {yelland-2015-new}.
  \end{inparaenum}
  Fatalmente, essas características essenciais de um \textit{prettyprinter}
  normalmente o restringem a uma linguagem específica ou, em alguns casos, a
  uma família de linguagens que compartilham semelhanças sintáticas ou
  estruturais.

  %%%%%%%%%%%%%%%%%%%%%%%%%%%
  %% The Tree-Sitter System %
  %%%%%%%%%%%%%%%%%%%%%%%%%%%

  \chapter{O Sistema \treesitter}
  \label{chapter:tree-sitter}

  O \treesitter é um sistema multilíngue de análise sintática para ferramentas
  de programação inicialmente desenvolvido como um projeto secundário por Max
  Brunsfeld. Como relatado pelo próprio autor \cite{github-2017-tree}, o
  \treesitter surgiu como uma tentativa de solucionar problemas presentes nas
  ferramentas de análise sintática da época. Mais especificamente, o sistema
  tinha como objetivos
  \begin{inparaenum}
    \item ser utilizado no ambiente de desenvolvimento para produzir árvores de
          sintaxe a partir da análise de códigos escritos em várias linguagens;
    \item implementar a análise incremental, permitindo a atualização da árvore
          de sintaxe em tempo real;
    \item expor através da árvore de sintaxe os nós representando suas
          construções gramaticais no código (e.g., classes, funções,
          declarações, etc.), diferentemente das ferramentas contemporâneas,
          que utilizavam uma abordagem simplística baseada em expressões
          regulares; e, por fim,
    \item ser livre de dependências, assim beneficiando sua adoção e
          aplicabilidade.
  \end{inparaenum}

  Além das funcionalidades previamente mencionadas presentes no \treesitter, o
  sistema também conta uma ferramenta de linha de comando que pode ser
  utilizada para gerar \textit{\gloss{parser}s} para uma linguagem a partir de
  sua gramática. A gramática é definida via a linguagem de programação
  JavaScript, a qual
  \begin{inparaenum}
    \item foi eleita 15 vezes seguidas pela
          \textit{Stack Overflow Developer Survey}
          \cite{stack-overflow-2022-stack} como a linguagem de programação mais
          comumente usada; além de
    \item ser amplamente considerada pela comunidade de programadores
          como uma das linguagens mais fáceis de aprender e programar
          \cites{berkeley-2023-11}{goel-2023-how}{w3schools-2023-javascript}.
  \end{inparaenum}
  A ferramenta de geração de \textit{parsers} também disponibiliza funções
  preestabelecidas para permitir a criação de gramáticas com diferentes níveis
  de complexidade. Não é surpreendente que, devido a essas características e
  facilidades presentes na criação de \textit{parsers}, exista uma variedade de
  linguagens de programação e formatos de arquivos -- variando de linguagens
  com sintaxes complexas, como C++ e Perl, a formatos de arquivos mais
  específicos, como \boxed{\texttt{.vhs}} e \boxed{\texttt{.rasi}} --, os quais
  têm parsers gerados pelo \treesitter e, consequentemente, são suportados pelo
  sistema.

  Na análise de código, é comum realizar tarefas que envolvem a busca de
  padrões em árvores sintáticas. Para isso, o \treesitter oferece uma pequena
  linguagem de consulta declarativa que é capaz de expressar esses padrões por
  meio de \textit{\gloss[Word]{sexp}s} e buscar correspondências. A linguagem
  de consulta suporta operadores que permitem
  \begin{inparaenum}
    \item a captura de nós;
    \item a quantificação de nós, análoga às expressões regulares;
    \item o agrupamento de nós;
    \item as alternâncias de nós;
    \item o uso de \textit{\gloss{wildcard}s}; e
    \item a ordenação de nós.
  \end{inparaenum}
  Adicionalmente, é possível definir propriedades nos nós da árvore sintática
  usando a linguagem de consulta do \treesitter. Também é permitido o uso de
  predicados -- funções arbitrárias para filtrar nós ou realizar verificações
  mais complexas durante a busca de padrões -- sejam eles
  \textit{\gloss{builtin}} ou estendidos por meio de uma API.

  Um exemplo da relevância e utilidade da linguagem de consulta é o plugin
  \textit{nvim-treesitter} \cite{nvim-treesitter-2023-nvim}, que é
  frequentemente utilizado no editor \textit{Neovim}
  \cite{neovim-2023-hyperextensible}. Esse plugin utiliza a linguagem de
  consulta para definir diferentes recursos, tais como
  \begin{inparaenum}
    \item \textit{\gloss{code-folding}}, que permite ocultar blocos de código;
    \item highlights, que realçam a sintaxe do código;
    \item indentações, que definem a estrutura do código;
    \item injeções, que permitem adicionar novas sintaxes a arquivos
          existentes; além de
    \item captura de nós correspondentes a construções gramaticais
          (e.g., funções, classes, métodos, etc.), os quais são
          frequentemente utilizados em rotinas de programação tais como a
          remoção e navegação do código.
  \end{inparaenum}

  Em resumo, o sistema \treesitter se mostra uma solução inovadora e eficiente
  para a análise sintática de códigos em diversas linguagens de programação,
  com uma abordagem diferenciada e sofisticada que possibilita a atualização em
  tempo real das árvores sintáticas e a identificação precisa das construções
  gramaticais presentes no código. Além disso, a ferramenta de linha de comando
  disponível no sistema facilita a geração de parsers a partir de gramáticas
  definidas em JavaScript, o que torna o processo mais acessível e
  personalizável para os programadores. Finalmente, a linguagem de consulta
  declarativa oferecida pelo \treesitter se mostra uma importante
  \textit{\gloss{feature}}, sendo utilizada em diversos plugins de editores de
  código para realizar tarefas variadas e sofisticadas, contribuindo
  significativamente no ambiente de desenvolvimento.

  %%%%%%%%%%%%%%%%%%%%%%
  %% The Witch Cooking %
  %%%%%%%%%%%%%%%%%%%%%%

  \chapter{O \witchcooking}

  \section{Motivações, Propósitos e Panorama}

  As motivações para o desenvolvimento do \witchcooking surgiram das
  dificuldades enfrentadas pelo autor ao estudar estilos de formatação de
  código-fonte em diversas linguagens de programação com as quais ele
  trabalhava. O estudo envolvia a exploração e análise tanto de estilos
  convencionais%
  \footnote
    {Neste trabalho, são considerados convencionais os estilos de formatação
     abordados nas páginas \textit{Indentation style} e
     \textit{Programming style} da Wikipédia.
     \cites{wiki-2023-indentation}{wiki-2023-programming}},
  dos quais muitos eram amplamente adotados, quanto de estilos
  não-convencionais. Durante a análise dos estilos convencionais de formatação,
  \begin{inparaenum}
    \item era imprescindível recorrer a diferentes \textit{prettyprinters},
    \item cada um com suas próprias configurações e níveis de suporte para
          esses estilos. Por outro lado, durante a análise dos estilos
          não-convencionais,
    \item a aplicação manual era inevitável, consumindo consideravelmente tempo
          e esforço.
  \end{inparaenum}

  Tendo em vista
  \begin{inparaenum}
    \item as dificuldades mencionadas anteriormente; além da
    \item complexidade envolvida no \textit{prettyprinting} --
    \item agravada pela proposta de formatação multilíngue --,
  \end{inparaenum}
  o \witchcooking foi desenvolvido como um software de linha de comando, de
  natureza prototípica, para a formatação de código-fonte. Para tal finalidade,
  a ferramenta requer a submissão do código-fonte a ser formatado, juntamente
  com o arquivo contendo as instruções de formatação. Em alinhamento com sua
  natureza experimental, o \witchcooking proporciona o código-fonte formatado
  resultante como saída, sem modificar o conteúdo original, efetivamente
  estabelecendo uma formatação segura.

  A ferramenta apresenta a seguinte usagem
  \boxed
    {\texttt{\textbf{cook}
     [\textbf{-l} \textit{LANG}]
     \textbf{-q} \textit{QUERY}
     [\textit{SRC}]}},
  onde
  \begin{inparaenum}
    \item \boxed{\texttt{\textbf{cook}}} é o nome do executável,
    \item \boxed{\texttt{[\textbf{-l} \textit{LANG}]}} é uma opção para
          explicitar a linguagem do código-fonte a ser formatado,
    \item \boxed{\texttt{\textbf{-q} \textit{QUERY}}} é utilizado para
          definir o caminho para o arquivo -- escrito na linguagem de consulta
          do \treesitter{} -- no qual as instruções de formatação estão
          definidas, e
    \item \boxed{\texttt{[\textit{SRC}]}} é um argumento opcional para
          especificar o caminho para o arquivo a ser formatado.
  \end{inparaenum}
  Quando \boxed{\texttt{[\textit{SRC}]}} é provido e
  \boxed{\texttt{[\textbf{-l} \textit{LANG}]}} é omitido, o software tentará
  inferir uma linguagem de acordo com a extenção do nome do arquivo. Por outro
  lado, quando \boxed{\texttt{[\textit{SRC}]}} é omitido, o código-fonte a ser
  formatado deve ser provido via \textit{\gloss{pipeline}}, e a sua linguagem
  deve ser explicitada via \boxed{\texttt{[\textbf{-l} \textit{LANG}]}}.

  \section{Instruindo a Formatação}

  Como mencionado na seção anterior, o \witchcooking formata o código-fonte
  conforme as instruções presentes no arquivo submetido via
  \boxed{\texttt{\textbf{-q} \textit{QUERY}}}, o qual é escrito na linguagem de
  consulta do \treesitter. Essas instruções constituem-se da
  \begin{inparaenum}
    \item busca de padrões -- i.e., \textit{pattern matching} -- em árvores
          sintáticas, da
    \item configuração, e da
    \item aplicação de predicados.
  \end{inparaenum}
  No entanto, como abordado no \autoref{chapter:tree-sitter}, a linguagem de
  consulta do \treesitter concebia originalmente a busca de padrões em árvores
  sintáticas. Portanto, para viabilizar a formatação de código-fonte, o
  \witchcooking introduz um conjunto adicional de configurações e
  predicados.

  %%%%%%%%%%%%%%%
  %% Conclusion %
  %%%%%%%%%%%%%%%

  \chapter{Conclusão}

  %%%%%%%%%%%%%%%%%%%%%%%%%
  % Post Textual Elements %
  %%%%%%%%%%%%%%%%%%%%%%%%%

  \postextual
  \printbibliography
  \printgloss{glossary}
\end{document}